% !TeX document-id = {ac384373-0975-4c4e-b6b2-65611dc82067}
% !TeX TXS-program:compile = txs:///pdflatex/[--shell-escape]
\documentclass[a4paper,12pt]{article}

\usepackage[T2A]{fontenc}
\usepackage[russian]{babel}
\usepackage[intlimits,namelimits]{amsmath}
\usepackage{amssymb}
\usepackage{mathtools}
\usepackage{mathrsfs}
\usepackage{geometry}
\usepackage[color={1 0 .5}]{attachfile2}

\usepackage{indentfirst}
\usepackage{icomma}
\usepackage{graphicx}
\usepackage{geometry}
\usepackage{setspace}
\usepackage{fancyhdr}
\usepackage{titlesec}
\usepackage{minted}
\usepackage{float}
\usepackage{caption}
\usepackage{rotating}

\usepackage{minted}

%%\usepackage[bibstyle=gost-numeric,citestyle=numeric-comp]{biblatex}

\usepackage{hyperref}

%%\addbibresource{biblio.bib}

\geometry{a4paper, nomarginpar, left=20mm, right=10mm, top=20mm, bottom=20mm}

\setlocalecaption{russian}{ref}{Список использованной литературы}


\newcommand\Vect[1]{{\boldsymbol{#1}}}
\newcommand{\dd}{\mathrm{d}}
\newcommand{\diag}{\operatorname{diag}}
\renewcommand{\Re}{\operatorname{Re}}
\renewcommand{\Im}{\operatorname{Im}}
\renewcommand{\exp}{\operatorname{e}}
\renewcommand{\imath}{\mathrm{i}}
\newcommand{\trp}{\mathrm{T}}

\newcommand{\mintedlabel}[1]{\phantomsection\label{\detokenize{#1}}}


% \DeclarePairedDelimeter\bra{\langle}{|}
\DeclarePairedDelimiter\ket{|}{\rangle}
\DeclarePairedDelimiter\abs{\lvert}{\rvert}
\DeclarePairedDelimiterX\commut[2]{[}{]}{#1,#2}
\DeclareCaptionLabelFormat{figRu}{Рисунок~#2}
\DeclareCaptionLabelSeparator{figRuSep}{~---~}
\captionsetup[figure]{labelformat=figRu,labelsep=figRuSep}

%\titleformat{\subsection}[block]{\filcenter}{\thesubsection}{0.5em}{}
\titleformat{\section}{\normalfont\Large\bfseries}{\thesection.}{1em}{}
\titlespacing{\section}{\parindent}{3.5ex plus 1ex minus .2ex}{2.3ex plus .2ex}
\titlespacing{\subsection}{\parindent}{3.5ex plus 1ex minus .2ex}{2.3ex plus .2ex}
\titlespacing{\subsubsection}{\parindent}{3.5ex plus 1ex minus .2ex}{2.3ex plus .2ex}

\fancypagestyle{plain}{%
	\fancyhf{}%
	\fancyfoot[R]{\thepage}%
	\renewcommand{\headrulewidth}{0pt}%
	\renewcommand{\footrulewidth}{0pt}%
}

%%%%%%%%%%%%%%%%%%%%%%%%%%%%%%%%%%%%%%%%%%%%%%%%%%%%%%%%%%%%%%%%%%%%%%%%%%%%%%%%

\makeatletter
\def\@sect[#1]#2{%
	\refstepcounter{app@section}
	\setcounter{equation}{0}
	\addcontentsline{toc}{section}{Приложение~\theapp@section}
	{%
		{
			\raggedleft%
			\normalfont%
			Приложение~\theapp@section\par%
		}%
		\Large\bfseries%
		#2%
		\par%
	}
	\vskip 3ex
}
% This is supposed to be "starred" version, but it doesn't have sense in the
% appendix, so we simply repeat above definition.
\def\@ssect#1{%
	\refstepcounter{app@section}
	\addcontentsline{toc}{section}{Приложение~\theapp@section}%
	{%
		{%
			\raggedleft%
			\normalfont%
			Приложение~\theapp@section\par%
		}%
		\Large\bfseries%
		#1%
		\par%
	}%
	\vskip 3ex
}

\renewcommand\appendix{%
	\newpage%
	\newcounter{app@section}
	\setcounter{app@section}{0}
	%\setcounter{equation}{0}%
	\renewcommand\theapp@section{\Asbuk{app@section}}
	\renewcommand\theequation{\theapp@section\arabic{equation}}
	%%% We need to redefine \section command here...
	%%% NB: this is messy, VERY messy (see GOST for autoreferat!)
	\renewcommand{\section}{%
		\secdef\@sect\@ssect
	}
}

%%%%%%%%%%%%%%%%%%%%%%%%%%%%%%%%%%%%%%%%%%%%%%%%%%%%%%%%%%%%%%%%%%%%%%%%%%%%%%%%

\allowdisplaybreaks[4]
\begin{document}

\onehalfspacing
\pagestyle{plain}

\begin{titlepage}
	\begin{center}%
		Министерство науки и высшего образования Российской Федерации\\
		федеральное государственное бюджетное образовательное учреждение\\
		высшего образования \\
		«Иркутский государственный университет»\\
		(ФГБОУ ВО «ИГУ»)\\
	\end{center}
	\vspace*{3em}
	
	\begin{flushright}
		\begin{minipage}{92mm}
			Физический факультет\\
			Кафедра теоретической физики\\
			Зав. кафедрой доцент, к.ф.-м.н.,\\
			Ловцов С.В.\;\underline{\hspace{35mm}}
		\end{minipage}
	\end{flushright}
	\vspace*{4em}
	
	\begin{center}
		\large
		\bfseries%
		ОТЧЕТ ПО ПРОИЗВОДСТВЕННОЙ ПРАКТИКЕ\\
		«Научно-исследовательская работа»\\[2ex]
		
		\Large
		Обработка данных для выявления качественной характеристики нейтринных осцилляций
	\end{center}
	
	\vspace*{3em}
	
	\begin{flushright}
		\begin{minipage}{95mm}
			Руководитель практики:
			\underline{\hspace{35mm}}\\
			к.ф.-м.н., доцент, Ломов В.П.\\[3ex]
			Студент гр. 01311-ДБ\\
			\underline{\hspace{35mm}}/Данеко Илья Игоревич\\[3ex]
			Работа защищена с оценкой \underline{\hspace{30mm}}\\
			<<\underline{\hspace{8mm}}>>\underline{\hspace{22mm}}\,2025г.\\
			Нормконтролёр\\
			\underline{\hspace{35mm}}/Фролова А.С./\\
		\end{minipage}
	\end{flushright}
	
	\vfill
	\begin{center}
		Иркутск 2025
	\end{center}
\end{titlepage}

\setcounter{page}{2}

\thispagestyle{empty}
\tableofcontents

\newpage

	\section*{Введение}
	\addcontentsline{toc}{section}{Введение}
Нейтринная физика сейчас является одним из рубежей современной физики. Также благодаря своим необычным свойствам нейтрино являются важным элементом для построения моделей за пределами стандартной модели.

Нейтрино --- это общее название нейтральных фундаментальных частиц с полуцелым спином, участвующих только в слабом и гравитационном взаимодействиях, и относящихся к классу лептонов. Наиболее важным открытием в физике нейтрино стало наблюдение осцилляций, которые к настоящему моменту подтверждены многочисленными экспериментами с солнечными, атмосферными, реакторными и ускорительными нейтрино.

Идея нейтринных осцилляций была предложена Бруно Понтекорво в 1957-1958 годах, она была основана по аналогии с K-мезонами \cite{11, 22}. Уже в 1962 году, Маки, Накагавой и Сакатой была разработана теория смешивания нейтрино в вакууме \cite{33}. В 1963 году Накагавой, Окониги, Сакатой и Тойодой была предложена связь между перемешиванием нейтрино с определенным флейвором и нейтрино с определённой массой \cite{44}. Вольфенштейном было обращено внимание на эффект при распространении нейтрино через обычное вещество (среда электронейтральная, немагнитная, нерелятивистская) в 1978 году \cite{55, 66}. В 1986 году Михеевым и Смирновым этот эффект был физически описан (MSW-эффект) \cite{77, 88}. И наконец, в 2015 году в качестве доказательства осцилляций нейтрино были признаны результаты SNO (Садберийская нейтринная обсерватория) \cite{99} и Super-Kamiokande \cite{1010}.

По сравнению с вакуумными осцилляциями нейтрино, при изучении осцилляции нейтрино в веществе, все сложнее. Поэтому в таком случае, как правило, применяют численные расчёты.

Основной целью данной работы стоит поиск качественной характеристики численного решения уравнения осцилляции нейтрино в среде.

В первом разделе мы кратко рассматриваем алгоритм обработки данных численного решения уравнения осцилляций в среде полученного с помощью метода DOPRI

	\newpage
	\section{Алгоритм для расчёта уравнения осцилляций нейтрино в среде}
	Поставленна задача найти качественные характеристики численного решения уравнения осцилляции нейтрино в среде
	\begin{equation}\label{eq:15}
		\imath\frac{\dd\Psi(\xi)}{\dd\xi}=[H_{0}+ v(\xi)W]\Psi(\xi),
	\end{equation}
	с начальным условием
	\begin{equation}\label{eq:19}
		\Psi(\xi_0)=
		\begin{pmatrix}
			c_{13}c_{12}\\
			s_{12}c_{13}\\
			s_{13}
		\end{pmatrix}.
	\end{equation}
	
	 Из численного анализа мы заметили, что реальная и мнимая часть \(\psi_1\) часто пересекают ось абсцисс, однако к правой границе области они выходит на асимптотическое поведение, более не пересекая ось абсцисс. Из подобных соображений было принято решение узнать как распределены нули \(\psi_{1}\). Для этого было введено понятие квазипериода. Квазипериод --- это длина отрезка содержащего в себе 3 точки в которых функция  (в нашем случае реальная или мнимая часть \(\psi_1\)) зануляется, причём 2 из них являются границами этого отрезка.

	Мы можем считать квазипериод функцией от точки \(\xi_0\) интервала \([0,1 ; 1]\) в таком смысле: находим 3 ближайших к \(\xi_0\) значения \(\xi\) при которых функция зависящая от \(\xi\) (в нашем случае реальная или мнимая часть \(\psi_1(\xi)\)) зануляется. Расстояние между крайними из них и есть искомый квазипериод в точке \(\xi_0\).

	Нас заинтересовало поведение квазипериода на всём интервале \([0,1 ; 1]\). Вполне логично построить график зависимости квазипериода от \(\xi\). Для получения данных по всему интервалу мы воспользовались тем, что на правом краю интервала \(\psi_{1}\) выходит на асимптотическое поведение более не пересекая ось абсцисс. Это означает, что есть самый крайний правый интервал с нулями и самый крайний квазипериод. Начиная с него и двигаясь влево мы получили необходимые данные. Для этого мы находим крайний правый квазипериод. Следующим этапом берём интервал последнего вычисленного квазипериода и смещаем его на \(\frac{2}{3}\) последнего вычисленного квазипериода влево. Далее ищем все зануления функции зависящей от \(\xi\) (в нашем случае реальной и мнимой части \(\psi_1(\xi)\)) внутри получившегося интервала и ищем квазипериод по приведённому выше алгоритму прировняв середину получившегося интервала к \(\xi_0\). Повторяем этапы начиная с получения нового интервала до тех пор пока не выполниться одно из двух условий. Либо в новом интервале окажется меньше 3 точек, либо количество вычисленных квазипериодов превысит заранее заданное число (мы выбрали 100). Это делается с помощью функции "funcQPeriodStat" в приведённом ниже коде для textsf{Mathematica}.

	Нам стало интересно насколько чаще реальная и мнимая часть \(\psi_{2,3}\) пересекают ось абсцисс по сравнению с реальной и мнимой частью \(\psi_1\) соответственно. Для этого мы посчитали количество занулений реальной и мнимой частей \(\psi_{2,3}\) на интервалах ранее посчитанных квазипериодов реальной и мнимой частей \(\psi_1\) соответственно. Это делается с помощью функции "funcCountZrsInRanges" в приведённом ниже коде для textsf{Mathematica}.
	
	Код textsf{Mathematica} для численного решения уравнения \eqref{eq:15} и обработки полученных данных:
	\inputminted[linenos,breaklines,fontsize=\scriptsize]{Markdown}{Sun_DOPRI_p5.md}

	\newpage
	\section*{Заключение}
	\addcontentsline{toc}{section}{Заключение}
	В работе рассматривался алгоритм для расчёта квазипериода реальной и мнимой части \(\psi_{1}\), а также количество занулений реальной и мнимой части \(\psi_{2,3}\) на интервале квазипериода реальной и мнимой части \(\psi_{1}\) при разных \(\xi\).

	Кроме этого отметим, что использование готовых решений для численного интегрирования, например, процедуры \texttt{NDSolve} программы \textsf{Mathematica} затрудняет воспроизводимость результатов, если только не задавать специально параметры для этой процедуры. К тому же, получаемый результат не представляет собой решение. Возвращаемый процедурой \texttt{NDSolve} результат — это интерполяционная формула. Но точных её параметры и другие характеристики неизвестны.

	Напоследок отметим, что для уменьшения риска человеческой ошибки были разработаны сценарии для  автоматизации большинства операций.


\newpage
\bibliography{biblio}
\bibliographystyle{ugost2008}
\addcontentsline{toc}{section}{Список использованной литературы}

\newpage
\appendix

\end{document}

